%% kapa: exam=1 round=27 year=2016 period=2 form=A subject=appraisal_law subject_ko=감정평가관계법규 seq=32 number=32
%% source: Q-Net 감정평가사 기출문제 (cst00309) artlSeq=5169402 file=extracted/1008/5169402/제27회 감정평가사 1차 문제/제27회 감정평가사 2교시문제(A형).pdf
%% conversion: kordoc 4.24.51 -> tools/gen1.py (auto); review=proofread-sonnet; flags=none
\begin{question}{32}
부동산등기법령상 소유권 등기에 관한 설명으로 옳은 것은?
\begin{choices}
  \item 등기관이 소유권보존등기를 할 때에는 등기원인과 그 연월일을 기록하지 아니한다.
  \item 토지대장에 최초의 소유자로 등록되어 있는 자의 상속인은 소유권보존등기를 신청할 수 없다.
  \item 등기관이 직권으로 소유권보존등기를 할 수 있는 경우는 없다.
  \item 소유권의 일부이전등기를 할 때 이전되는 지분을 표시하지 않아도 된다.
  \item 소유권의 이전에 관한 사항은 등기기록의 을구에 기록한다.
\end{choices}
\end{question}
\answer{1}
